\documentclass[11pt]{article}
\usepackage[T1]{fontenc}
\usepackage{lmodern}
\usepackage[margin=1in]{geometry}
\usepackage{amsmath,amssymb,booktabs,array,enumitem}
\usepackage{pdfpages}
\usepackage[colorlinks=true,linkcolor=blue!55!black,urlcolor=blue!55!black]{hyperref}
\usepackage{microtype}
\setlength{\parindent}{0pt}
\setlength{\parskip}{0.35em}
\hypersetup{pdftitle={Wave-Constrained Computation: Formal Equivalence, Physical Resource Bounds, and the P versus NP Problem},pdfauthor={Richard J. Reyes},pdfsubject={Consolidated PVNP technical volume}}
\title{\bfseries Wave-Constrained Computation\\[0.4em]\large Formal Equivalence, Physical Resource Bounds,\\and the P versus NP Problem\\[0.8em]\normalsize Consolidated Technical Volume}
\author{Richard J. Reyes}
\date{Consolidated edition: November 2026\\\small Original studies released in 2025}
\begin{document}
\maketitle
\begin{abstract}
This technical volume binds the three Wave-Constrained Computation (WCC) manuscripts and an independent fixed-mass nodal-coercivity supplement into one navigable PDF. The earlier revised documents are reproduced in full as source-preserving parts. This introduction distinguishes exact uniform discrete WCC results, charged physical wave realizations, restricted lower bounds, and unresolved universal separation claims. Neither classical P versus NP nor physical WCC separation is proved by this collection.
\end{abstract}
\tableofcontents
\clearpage

\section{Purpose and publication provenance}
This is a consolidated edition of earlier research, not four additional newly authored papers. The historical titles, mathematical qualifications, original citations, and internal references are preserved in the bound originals. A later corrected theorem takes precedence over a broader historical informal formulation.

\begin{center}\small
\begin{tabular}{@{}p{0.13\linewidth}p{0.50\linewidth}p{0.29\linewidth}@{}}\toprule
Part & Original work & Historical publication\\\midrule
I & Curvature-bounded wave computation, May 7, 2025 & \href{https://doi.org/10.5281/zenodo.17743607}{Zenodo 17743607}\\
II & Discrete WCC and Turing equivalence, November 26, 2025 & \href{https://doi.org/10.5281/zenodo.17732642}{Zenodo 17732642}\\
III & Physical complexity reassessment, December 1, 2025 & \href{https://doi.org/10.5281/zenodo.17783074}{Zenodo 17783074}\\
Appendix & Fixed-mass nodal coercivity, October 2026 & Independent auxiliary note\\
\bottomrule
\end{tabular}
\end{center}
These original-release dates are not reset by the 2026 LaTeX revisions or this compilation. The DOI links identify historical records, not a newly assigned consolidated-volume DOI.

\section{Results and scope}
\subsection{Uniform discrete computation}
With finite alphabets and local update rules, polynomial-time constructible lattices, and charged uniform encoders and decoders, the discrete model has the formally established class equivalences
\[
\mathbf P_{\mathrm{WCC}}^{\mathrm{disc}}=\mathbf P_{\mathrm{TM}},
\qquad
\mathbf{NP}_{\mathrm{WCC}}^{\mathrm{disc}}=\mathbf{NP}_{\mathrm{TM}}.
\]
The original two-phase machine \(W_M\) has a distinct open uniformity obligation: its complete original initial encoder and lattice writer must be certified polynomial-time transducers. The generic discrete class equivalence does not by itself discharge that obligation.

\subsection{Physical realization and mathematical lower bounds}
The continuum WCT resource model additionally requires uniform field preparation, precision, stable readout, and full accounting of energy and computational cost. The discrete equivalences do not prove that a concrete nonlinear PDE realizes them with polynomial overhead. Finite-domain coercivity, phase readout, and separator estimates have explicit assumptions that need not be input-size uniform.

The repository includes uniform Cook--Levin reductions and a characterization of the target dense-SAT lower bound. Restricted one-tape crossing-sequence arguments yield quadratic-type lower bounds. They do not show that every deterministic SAT algorithm requires superpolynomial time. In particular,
\[
\mathrm{SAT}_{\mathrm{dense}}\notin\mathbf P_{\mathrm{TM}}
\]
remains an open conjectural target. The physical all-decider separator and expander claims also remain open.

\section{Maintained formalization ledger and verification}
The amended 91-claim ledger records 78 Exact, 2 Restricted, 4 Refuted, and 7 Open claims. The Exact category includes proved conditional implications whose premises are not established; it does not mean classical P is unequal to NP. See the continuously maintained \href{https://github.com/rickyjreyes/PVNP/blob/main/pvnp/proofs/manuscript-lean-inventory.md}{manuscript formalization inventory} and \href{https://github.com/rickyjreyes/PVNP/tree/main/lean}{Lean source tree}.

\begin{center}\small
\begin{tabular}{@{}p{0.33\linewidth}p{0.58\linewidth}@{}}\toprule
Topic & Established limit\\\midrule
Uniform discrete WCC & Proved classical P and NP class equivalences within the formal discrete model.\\
Original two-phase \(W_M\) & Exact decoder and simulation support; full original polynomial-time writer obligations remain restricted.\\
Dense SAT & Uniform Cook--Levin and restricted quadratic single-tape lower bounds; universal superpolynomial target open.\\
Physical WCC & Charged, robust, uniform physical-to-discrete resource bridge missing.\\
Curvature interfaces & Restricted finite mesh and graph estimates; universal separator and expander hardness open.\\
Phase readout & Logical-energy SAT certification does not imply efficient physical convergence or preparation.\\
Auxiliary nodal estimate & Fixed-domain analytical result; its finite two-site algebra is formalized, not the full continuum theorem.\\
\bottomrule\end{tabular}
\end{center}

\textbf{Build caveat.} Some newly merged Phase XII and XIII Lean declarations entered the repository while GitHub-hosted checks failed before their first job step. This volume does not certify those declarations as independently kernel-checked until a complete Lean build and axiom audit succeeds. A PDF compilation is not evidence of a successful Lean verification.

\section{Reproduction and source maintenance}
The authoritative historical LaTeX sources remain under the four original manuscript directories in the repository. The single-release build first compiles them into temporary PDFs, binds them behind this overview, and installs only one primary consolidated PDF. Run \texttt{bash scripts/build\_pdfs.sh} to regenerate it. For the Lean formalization, use the pinned toolchain and run \texttt{lake build PVNP} followed by the standard axiom audit from within the \texttt{lean} directory.

The reproduced parts retain their own title pages, internal pagination, bibliographies, and reference numbering, to preserve provenance and avoid silently rewriting earlier published work. They constitute one \emph{bound research volume}, not yet an editorially merged single-paper text.

\clearpage
\phantomsection\addcontentsline{toc}{section}{Part I: Curvature-bounded wave computation}
\section*{Part I: Curvature-Bounded Wave Computation}
\emph{Original research May 7, 2025. Revised source reproduced in full.}
\clearpage
\includepdf[pages=-,pagecommand={},scale=0.97]{archive/part-i.pdf}

\clearpage
\phantomsection\addcontentsline{toc}{section}{Part II: Discrete WCC and Turing equivalence}
\section*{Part II: Discrete WCC and Turing Equivalence}
\emph{Original research November 26, 2025. Revised source reproduced in full.}
\clearpage
\includepdf[pages=-,pagecommand={},scale=0.97]{archive/part-ii.pdf}

\clearpage
\phantomsection\addcontentsline{toc}{section}{Part III: Physical complexity reassessment}
\section*{Part III: Physical Complexity Reassessment}
\emph{Original research December 1, 2025. Revised source reproduced in full.}
\clearpage
\includepdf[pages=-,pagecommand={},scale=0.97]{archive/part-iii.pdf}

\clearpage
\phantomsection\addcontentsline{toc}{section}{Auxiliary appendix: Fixed-mass nodal coercivity}
\section*{Auxiliary Appendix: Fixed-Mass Nodal Coercivity}
\emph{Independent auxiliary note, not one of the original 91 manuscript claims.}
\clearpage
\includepdf[pages=-,pagecommand={},scale=0.97]{archive/supplement.pdf}
\end{document}
